%-*-latex-*-
%-*-latex-*-
\newcommand\COURSE{ciss240}
\newcommand\ASSESSMENT{q1005}
\newcommand\ASSESSMENTTYPE{Quiz}
\newcommand\POINTS{\textwhite{xxx/xxx}}

\input{myquizpreamble}
\input{yliow}
\input{\COURSE}
%-*-latex-*-

%-*-latex-*-

%-*-latex-*-
\renewcommand\TITLE{\ASSESSMENTTYPE \ \ASSESSMENT}

\renewcommand\EMAIL{}
\renewcommand\AUTHOR{eagraham1@cougars.ccis.edu}

\textwidth=6in
\begin{document}
\topmatterthree

This is a closed-book, no compiler, 5 minute quiz.

%------------------------------------------------------------------------------
\nextq
Complete the following code fragment to compute
the sum $1 + 2 + 3 + \cdots + 1000$.
\\
\ANSWER
\begin{answercode}
      sum = 0; // declare and initialize
for (int i = 1; i <= 1000; i++)
{
    sum += i; 
}
std::cout << sum << std::endl;
\end{answercode}

%------------------------------------------------------------------------------
\nextq
Complete the following code fragment to compute
the sum $10 + 12 + 14 + \cdots + 996 + 998 + 1000$.
\\
\ANSWER
\begin{answercode}
      sum = 0; // declare and initialize
for (int i = 10; i <= 1000; i += 2)
{
    sum += i;
}
std::cout << sum << std::endl;
\end{answercode}

%------------------------------------------------------------------------------
\nextq
Complete the following code fragment to compute
the sum of the integers 10, 12, 14, ..., 996, 998, 1000
\textit{except} for the integers
in [100, 200], i.e., compute and print
\[
10 + 12 + 14 + \cdots + 94 + 96 + 98 + 202 + 204 + 206 + \cdots + 996 + 998 + 1000
\]
You must use only one for-loop.
Of course this can be done in two separate for-loops, but for this
question you have to following the given code fragment
and use only one for-loop.
(You need an if statement.)
\\
\ANSWER
\begin{answercode}
      sum = 0; // declare and initialize
for (int i = 10; i < 1000; i += 2)
{
    if(i != 100 || i != 200)
    {
        sum += i; 
    }
}
std::cout << sum << std::endl;
\end{answercode}

%------------------------------------------------------------------------------
%-*-latex-*-
\newpage
\input{instructions}
\end{document}

